\PassOptionsToPackage{numbers}{natbib}
\documentclass{article}
\usepackage{iclr2024_conference}
\usepackage{luatexja-fontspec}
\usepackage{hyperref}
\usepackage{url}
\usepackage{booktabs}
\usepackage{amsmath}
\usepackage{amsfonts}
\usepackage{graphicx}
\graphicspath{{../}}

\InputIfFileExists{values.tex}{}{}
\providecommand{\airasval}[1]{\textbf{??airasval:\detokenize{#1}??}}
\providecommand{\unverified}[1]{#1}
\providecommand{\airasrecordlink}[1]{#1}

\title{SciGym における 2026 年 9 月時点の最新モデル 6 種の事前登録評価: 論文 Table 1 と同じ条件で small を測り直し、large にも広げる}
\author{AIRAS}

\begin{document}
\maketitle

\begin{abstract}
SciGym \citep{duan-2025-measuring} は、反応をすべて取り除いた SBML モデルを LLM エージェントに渡し、摂動実験とコード実行を繰り返して反応機構を復元させるベンチマークである。論文の Table 1 は 2025 年春の 3 社のフロンティアモデル 6 つを SciGym-small 137 件で各 1 回走らせた平均値を報告している。本研究は、公開コードの Controller と評価器をそのまま用い、論文と同じ条件（137 件、最大 20 反復、temperature 1.0）で、2026 年 9 月時点で Vercel AI Gateway から利用できる各社の最新モデル 6 種（OpenAI の gpt-6-sol と gpt-6-luna、Anthropic の claude-opus-5.5 と claude-sonnet-5、Google の gemini-3.1-pro-preview と gemini-3.8-flash）を評価し、論文の Table 1 をそのまま並べて報告する。あわせて、論文が評価していない large split 213 件も同じ条件で走らせる。実験前に、small と large のそれぞれで 6 モデルすべての平均 Simulation Trajectory Error（STE）が論文 Table 1 の最良行（Gemini-2.5-Pro の 0.3212）以下であることを反証線として登録する。Reaction Matching Score（RMS）と Network Topology Score（NTS）は同じ実行から表として報告する。
\end{abstract}

\section{はじめに}
\label{sec:intro}
LLM に科学的発見のサイクル全体（実験の設計、結果の解析、機構の提案）を行わせて測るベンチマークとして、SciGym \citep{duan-2025-measuring} が提案された。BioModels 由来の SBML モデルから反応をすべて取り除いたものをエージェントに渡し、エージェントは初期濃度を変えた摂動実験を要求し、返ってきた時系列を Python で解析し、最終的な SBML を提出する。提出モデルは真のモデルと、軌道誤差（STE）と反応の一致度（RMS）で比較される。

論文の主結果である Table 1 は、3 社の pro / mini 計 6 モデルを SciGym-small の 137 件で評価した平均値である \cite[s1.p2, s1.p10]{duan-2025-measuring}。評価されたのは 2025 年春の版（Claude-3.7-Sonnet と 3.5-Haiku、Gemini-2.5-Pro と Flash のプレビュー版、GPT-4.1 と 4.1-mini）である \cite[s1.p1]{duan-2025-measuring}。その後 1 年半で各社のモデルは数世代進んだが、同じ閉ループの課題がどこまで解けるようになったかは示されていない。

AIRAS の先行研究は、現行 API の Gemini-2.5-Pro で同じ課題を解かせたときの平均 STE が論文の 0.3212 から 0.05 以内に収まるかを事前登録して検証し \cite[s3.p1]{scigymtable1reproduction-2026-scigym}、また open-weight モデル 6 種が論文の GPT-4.1-mini の水準（0.6007）に届くかを同じプロトコルで評価した \cite[s4.p1]{scigymrikyuopenmodels-2026-scigym}。後者では kimi-k2.6 の平均 STE が \unverified{0.2100} \cite[s4.p2]{scigymrikyuopenmodels-2026-scigym}、kimi-k3 が \unverified{0.2251} \cite[s4.p3]{scigymrikyuopenmodels-2026-scigym} と、論文 Table 1 の最良行 Gemini-2.5-Pro の \unverified{0.3212} \cite[s1.p5]{duan-2025-measuring} を下回った。本研究は同じ Controller、同じ条件で 2026 年 9 月時点の各社最新 6 モデルを走らせ、論文の Table 1 をそのまま並べて世代の差を報告する。論文が計算資源の制約から評価しなかった large split 213 件 \cite[s1.p2]{duan-2025-measuring} にも同じ設定で広げる。

判定は AIRAS の事前登録の枠組みに従い、仮説・判定基準・予測区間・実行条件を実験前に記録へ凍結し、実験後は計測値だけを流し込む。

\section{関連研究}
\label{sec:related}
SciGym \citep{duan-2025-measuring} は、単一スキルを測る既存のベンチマークと異なり、エージェントが自ら摂動実験を設計してその結果を受け取る閉ループを評価する。評価指標は 3 つで、種間の相互作用の F1（NTS）、反応物と生成物の集合が一致した反応の適合率・再現率・F1（RMS。modifier まで要求する厳格版もある）\cite[s1.p8]{duan-2025-measuring}、および全種の時系列の対称平均絶対パーセント誤差（STE）\cite[s1.p9]{duan-2025-measuring} である。論文は pro モデルが mini モデルを一貫して上回り \cite[s1.p6]{duan-2025-measuring}、Gemini-Pro が最良であった \cite[s1.p7]{duan-2025-measuring} と報告している。

先行する Table 1 の再現では、現行 API の GPT-4.1 と 4.1-mini、Gemini-2.5-Flash は論文値の 0.05 以内に収まったが、Gemini-2.5-Pro は \unverified{0.4148} \cite[s3.p2]{scigymtable1reproduction-2026-scigym} と論文値を再現しなかった。open-weight モデルの評価では kimi-k2.6 が \unverified{0.2100} \cite[s4.p2]{scigymrikyuopenmodels-2026-scigym}、kimi-k3 が \unverified{0.2251} \cite[s4.p3]{scigymrikyuopenmodels-2026-scigym}、deepseek-v4.1-flash が \unverified{0.2405} \cite[s4.p4]{scigymrikyuopenmodels-2026-scigym} に達した。

\section{仮説と予測}
\label{sec:hypothesis}
以下は記録（record.json）から生成された仮説・主張・判定基準・予測区間の一覧である。実験後は Observed と Verdict が計測値から埋められる。

% AUTO-GENERATED by the update_record tool. DO NOT EDIT.
% verify_paper regenerates this file from record.json and the run outputs
% and fails on any difference, so manual edits always surface.
\noindent\textbf{Hypothesis 1.} 2026 年 9 月時点で Vercel AI Gateway から利用できる各社の最新モデル 6 種（OpenAI の gpt-6-sol と gpt-6-luna、Anthropic の claude-opus-5.5 と claude-sonnet-5、Google の gemini-3.1-pro-preview と gemini-3.8-flash）は、SciGym 公開コードの Controller をそのまま用いて論文と同じ条件（SciGym-small 137 件、最大 20 反復、temperature 1.0）で解かせたとき、いずれも平均 Simulation Trajectory Error（STE）が論文 Table 1 の最良行（Gemini-2.5-Pro の 0.3212）以下である。~\cite[s1.p1, s1.p5, s1.p10]{duan-2025-measuring}, \cite[s3.p2]{scigymtable1reproduction-2026-scigym}, \cite[s4.p2]{scigymrikyuopenmodels-2026-scigym}
\begin{enumerate}
\item[\textbf{Claim 1}] gpt-6-sol（openai/gpt-6-sol）で SciGym-small 137 件を最大 20 反復で解かせたときの平均 STE は、論文 Table 1 の最良行（Gemini-2.5-Pro）の 0.3212 以下である。~\cite[s1.p2, s1.p3, s1.p4, s1.p5, s1.p10]{duan-2025-measuring}, \cite[s4.p2]{scigymrikyuopenmodels-2026-scigym}, \cite[s2.p1]{h4duan-2025-scigym}
  \emph{Rationale:} OpenAI の現行世代の推論モデルが論文の最良行に届けば、h1 の gpt-6-sol の部分を担う。
  \emph{Criterion:} \texttt{\detokenize{gpt-6-sol}}.\texttt{\detokenize{trajectory_smape}} $-$ 0.3212 $\leq 0$.
  \emph{Prediction:} $[-0.2, 0]$ (open-weight の kimi-k2.6 / kimi-k3 / deepseek-v4.1-flash が同条件で 0.21〜0.24 に達した（s4.p2〜s4.p4）。商用の最新モデルはそれと同等かそれ以上を予測する).
  \emph{Observed:} pending.
  \emph{Verdict:} pending.
\item[\textbf{Claim 2}] gpt-6-luna（openai/gpt-6-luna）で SciGym-small 137 件を最大 20 反復で解かせたときの平均 STE は、論文 Table 1 の最良行（Gemini-2.5-Pro）の 0.3212 以下である。~\cite[s1.p2, s1.p3, s1.p4, s1.p5, s1.p10]{duan-2025-measuring}, \cite[s4.p5]{scigymrikyuopenmodels-2026-scigym}, \cite[s2.p1]{h4duan-2025-scigym}
  \emph{Rationale:} OpenAI の最小クラスのモデルが論文の最良行に届けば、h1 の gpt-6-luna の部分を担う。
  \emph{Criterion:} \texttt{\detokenize{gpt-6-luna}}.\texttt{\detokenize{trajectory_smape}} $-$ 0.3212 $\leq 0$.
  \emph{Prediction:} $[-0.15, 0.05]$ (open-weight の 27B〜35B 級（qwen3.8-27b 0.25、qwen3.6-35b 0.46）が同条件で最良行の前後に分かれた（s4.p5、s4.p6）。最小クラスの商用モデルはその中間を予測する).
  \emph{Observed:} pending.
  \emph{Verdict:} pending.
\item[\textbf{Claim 3}] claude-opus-5.5（anthropic/claude-opus-5.5）で SciGym-small 137 件を最大 20 反復で解かせたときの平均 STE は、論文 Table 1 の最良行（Gemini-2.5-Pro）の 0.3212 以下である。~\cite[s1.p2, s1.p3, s1.p4, s1.p5, s1.p10]{duan-2025-measuring}, \cite[s4.p3]{scigymrikyuopenmodels-2026-scigym}, \cite[s2.p1]{h4duan-2025-scigym}
  \emph{Rationale:} Anthropic の現行の上位モデルが論文の最良行に届けば、h1 の claude-opus-5.5 の部分を担う。
  \emph{Criterion:} \texttt{\detokenize{claude-opus-5.5}}.\texttt{\detokenize{trajectory_smape}} $-$ 0.3212 $\leq 0$.
  \emph{Prediction:} $[-0.2, 0]$ (c1 と同じ（s4.p2〜s4.p4）).
  \emph{Observed:} pending.
  \emph{Verdict:} pending.
\item[\textbf{Claim 4}] claude-sonnet-5（anthropic/claude-sonnet-5）で SciGym-small 137 件を最大 20 反復で解かせたときの平均 STE は、論文 Table 1 の最良行（Gemini-2.5-Pro）の 0.3212 以下である。~\cite[s1.p2, s1.p3, s1.p4, s1.p5, s1.p10]{duan-2025-measuring}, \cite[s4.p4]{scigymrikyuopenmodels-2026-scigym}, \cite[s2.p1]{h4duan-2025-scigym}
  \emph{Rationale:} Anthropic の中位モデルが論文の最良行に届けば、h1 の claude-sonnet-5 の部分を担う。
  \emph{Criterion:} \texttt{\detokenize{claude-sonnet-5}}.\texttt{\detokenize{trajectory_smape}} $-$ 0.3212 $\leq 0$.
  \emph{Prediction:} $[-0.2, 0]$ (中位の商用モデルは open-weight の上位（0.21〜0.25）に近いと予測する（s4.p2〜s4.p5）).
  \emph{Observed:} pending.
  \emph{Verdict:} pending.
\item[\textbf{Claim 5}] gemini-3.1-pro-preview（google/gemini-3.1-pro-preview）で SciGym-small 137 件を最大 20 反復で解かせたときの平均 STE は、論文 Table 1 の最良行（Gemini-2.5-Pro）の 0.3212 以下である。~\cite[s1.p2, s1.p3, s1.p4, s1.p5, s1.p7, s1.p10]{duan-2025-measuring}, \cite[s3.p2]{scigymtable1reproduction-2026-scigym}, \cite[s2.p1]{h4duan-2025-scigym}
  \emph{Rationale:} 論文で最良だった Gemini Pro 系の現行版が自らの旧版の値に届けば、h1 の gemini-3.1-pro-preview の部分を担う。
  \emph{Criterion:} \texttt{\detokenize{gemini-3.1-pro}}.\texttt{\detokenize{trajectory_smape}} $-$ 0.3212 $\leq 0$.
  \emph{Prediction:} $[-0.2, 0]$ (c1 と同じ。なお現行 API の gemini-2.5-pro の再実行は 0.4148 で論文値を再現しなかった（s3.p2）ため、旧版との差は環境差を含む).
  \emph{Observed:} pending.
  \emph{Verdict:} pending.
\item[\textbf{Claim 6}] gemini-3.8-flash（google/gemini-3.8-flash）で SciGym-small 137 件を最大 20 反復で解かせたときの平均 STE は、論文 Table 1 の最良行（Gemini-2.5-Pro）の 0.3212 以下である。~\cite[s1.p2, s1.p3, s1.p4, s1.p5, s1.p10]{duan-2025-measuring}, \cite[s3.p3]{scigymtable1reproduction-2026-scigym}, \cite[s2.p1]{h4duan-2025-scigym}
  \emph{Rationale:} Gemini Flash 系の現行版が論文の最良行に届けば、h1 の gemini-3.8-flash の部分を担う。
  \emph{Criterion:} \texttt{\detokenize{gemini-3.8-flash}}.\texttt{\detokenize{trajectory_smape}} $-$ 0.3212 $\leq 0$.
  \emph{Prediction:} $[-0.15, 0.05]$ (c2 と同じ。現行 API の gemini-2.5-flash の再実行は 0.4465 だった（s3.p3）が、世代が進んだ分を見込む).
  \emph{Observed:} pending.
  \emph{Verdict:} pending.
\end{enumerate}
\noindent\emph{Assumed, so that the claims together imply Hypothesis 1:}
\begin{itemize}
\item temperature 1.0 での 1 回実行の平均 STE（137 件平均）の実行間変動が判定に対して十分小さいこと。先行する Table 1 再現では件ごとの STE の標準偏差はおよそ 0.27、137 件平均の標準誤差はおよそ 0.02 であった（c1〜c6）
\item Vercel AI Gateway 経由の各モデルが、提供元の API から直接呼んだ場合と同じ能力で応答すること。reasoning の強さは各モデルの既定値のまま変えない（c1〜c6）
\item 論文の 6 モデルは 2025 年春の版であり、本研究の 6 モデルは 2026 年 9 月時点の各社の最新版であること。世代の差を測るのが目的なので、同名モデルの再現は求めない（c1〜c6）
\item aarch64 向けの libroadrunner 2.7.0 と antimony 2.14.0 が、論文環境の libroadrunner 2.8.0 と同じ軌道と評価値を与えること（c1〜c6）
\item 平均 STE が「論文の最良行に届くか」を代表すること。RMS（modifier あり / なし）と NTS は同じ実行から表として報告するが、判定基準には用いない（c1〜c6）
\item ゲートウェイの予算超過（HTTP 402）で run が途中で止まらないこと。止まった場合はその run を最初からやり直す（c1〜c6）
\end{itemize}


\paragraph{判定基準と予測区間の根拠.}
仮説 1（small）の各主張の反証線は「（当該モデルの平均 STE）$-$ \unverified{0.3212} $\le 0$」である。\unverified{0.3212} は論文 Table 1 の最良行 Gemini-2.5-Pro の STE \cite[s1.p5]{duan-2025-measuring} であり、6 行すべてに同じ基準を置く。1 年半後の各社の最新モデルが、論文時点の最良に届くかどうかを問う。仮説 2（large）も同じ基準を置く。large には論文の公表値が無いので、small での最良行を絶対水準の基準とし、予測区間は上位モデルで $[-0.10, 0.10]$、最小クラスで $[-0.05, 0.15]$ と幅を広く取る。予測区間は上位モデル（gpt-6-sol、claude-opus-5.5、claude-sonnet-5、gemini-3.1-pro-preview）で $[-0.20, 0]$、最小クラス（gpt-6-luna、gemini-3.8-flash）で $[-0.15, 0.05]$ とする。根拠は open-weight の上位モデルが同条件で kimi-k2.6 \unverified{0.2100} \cite[s4.p2]{scigymrikyuopenmodels-2026-scigym}、kimi-k3 \unverified{0.2251} \cite[s4.p3]{scigymrikyuopenmodels-2026-scigym}、deepseek-v4.1-flash \unverified{0.2405} \cite[s4.p4]{scigymrikyuopenmodels-2026-scigym} に達したことと、27B〜35B 級が 0.25 と 0.46 に分かれたこと \cite[s4.p5, s4.p6]{scigymrikyuopenmodels-2026-scigym} である。1 件の STE の標準偏差を 0.27 程度と見積もると 137 件平均の標準誤差は約 0.02 で、差が 0.02 未満のときは結論が揺らぎうることを前提として記録する。

RMS（modifier あり / なし）と NTS も同じ実行から表として報告するが、1 つの主張に置ける判定基準は 1 つであるため、判定には用いない。各社の上位と下位の差、small と large の難易度差も表で示すにとどめ、主張にはしない。

\section{方法}
\label{sec:method}
\paragraph{タスク.}
各インスタンスは真の SBML（\texttt{truth.xml}）、反応を取り除いた不完全な SBML（\texttt{partial.xml}）、シミュレーション条件（\texttt{truth.sedml}）からなる。エージェントは不完全な SBML を受け取り、各ターンで実験（初期濃度の変更を指定し、真のモデルの時系列を得る）、コード実行（numpy / pandas / scipy / libsbml と、仮説モデルを走らせる \texttt{simulate} 関数が使える）、提出のいずれかを markdown 形式で返す。提出は任意の時点で行え、反復数の上限は 20 である \cite[s1.p3]{duan-2025-measuring}。提出した SBML が無効な場合は追加のデバッグ反復が許され、それでも無効なら不完全 SBML のまま採点される \cite[s1.p4]{duan-2025-measuring}。

\paragraph{実行系.}
ループ、プロンプト、実験の実行はすべて公開コード \citep{h4duan-2025-scigym} の \texttt{Controller} をそのまま用いる。本研究が書くのは、(1) OpenAI 互換 API でモデルを呼ぶ \texttt{LLM} 実装（公式の GPT 実装と同じく、システムプロンプトと全履歴を毎回送る）、(2) インスタンスごとに 1 プロセスで Controller を起動し、137 件を並列に走らせる駆動部、(3) 各インスタンスの提出モデルと公式データを評価層の入力ファイルに書き出す部分だけである。公式コードでは履歴がクラス属性として定義されており 1 プロセスで複数件を走らせると履歴が持ち越されるが、1 件 1 プロセスとすることでこの問題を避ける。

\paragraph{指標.}
採点は評価層 airas-eval のベンチマークパック \texttt{scigym\_small}（137 件）と \texttt{scigym\_large}（213 件）が行う。パックは公式リリースの各件（truth / partial / sedml）を sha256 で固定し、公開コードの \texttt{Evaluator}（コミット 88a7b93）で提出モデルを採点して、STE（\texttt{trajectory\_smape}）、RMS の適合率・再現率・F1（modifier なし \texttt{reaction\_*}、あり \texttt{reaction\_*\_with\_modifiers}）、NTS の F1（\texttt{topology\_f1}）を件の単純平均で返す。提出 SBML が無い、または無効な件は、公式評価どおり不完全 SBML の採点値で平均に含める。

\section{実験設計}
\label{sec:design}
\paragraph{モデル.}
Vercel AI Gateway（OpenAI 互換）経由で、2026 年 9 月時点の各社の最新モデルを上位と下位の 2 つずつ用いる: openai/gpt-6-sol と openai/gpt-6-luna、anthropic/claude-opus-5.5 と anthropic/claude-sonnet-5、google/gemini-3.1-pro-preview と google/gemini-3.8-flash。論文の pro / mini の構図に対応させるが、各社の製品構成に厳密な対応は無く、料金と提供元の位置づけから選んだ。reasoning の強さは各モデルの既定値のまま変えない。

\paragraph{データ.}
SciGym の公式リリース（Hugging Face \texttt{h4duan/scigym-sbml}）のうち、論文が評価対象とした small split 137 件 \cite[s1.p2]{duan-2025-measuring} に加え、論文では評価されていない large split 213 件を用いる（いずれも RIKYU 上のコピー）。large の真のモデルの SBML は small のおよそ 3.5 倍長い。評価層は各件の内容を公式リリースの sha256 と照合する。1 run は 1 モデル × 1 split で、計 12 run である。

\paragraph{設定.}
\texttt{max\_iterations} は論文どおり 20 \cite[s1.p3]{duan-2025-measuring}（公開 config の既定値は 5 \cite[s2.p2]{h4duan-2025-scigym}）、\texttt{eval\_debug\_rounds} は論文本文の 3 \cite[s1.p4]{duan-2025-measuring}（公開 config の既定値は 5 \cite[s2.p3]{h4duan-2025-scigym}）、\texttt{temperature} は公開 config の 1.0 \cite[s2.p1]{h4duan-2025-scigym}、摂動は初期濃度の変更のみとする。応答のトークン上限は 32768 とし、reasoning で使い切った応答は上限を倍にして呼び直す。本文が空で返った呼び出しは呼び直し、reasoning 側に本文が入っていればそれを用いる。各モデル、各インスタンスを 1 回ずつ実行する。ゲートウェイの予算超過（HTTP 402）が起きた run は途中の件を残さず最初からやり直す。

\paragraph{計算資源と環境.}
実行は Seyval を通じて RIKYU（Slurm、GB200 ノード、aarch64）で行い、LLM は Vercel AI Gateway を呼ぶ。GPU は使わないが、資源の単位が GPU であるため 1 GPU を要求する。環境は Dockerfile と \texttt{uv.lock} で固定する。aarch64 向けの wheel が無いため、libroadrunner は公式の 2.8.0 ではなく 2.7.0、antimony は 2.14.0 を用い、rrplugins と phrasedml は除外する。この差でシミュレーション結果がわずかに変わる可能性があり、仮説の前提として記録する。

\paragraph{段階.}
sanity は各 split の先頭 1 件を 2 反復で走らせて配管を確かめる。pilot は sorted した件から等間隔に 10 件を 20 反復で走らせ、1 件あたりの所要時間と費用を測る。主張の判定は full（small 137 件、large 213 件）の結果だけで行い、sanity と pilot の結果は記録に取り込まない。

\paragraph{run と主張の対応.}
small: \texttt{gpt-6-sol} $\to$ Claim 1、\texttt{gpt-6-luna} $\to$ Claim 2、\texttt{claude-opus-5.5} $\to$ Claim 3、\texttt{claude-sonnet-5} $\to$ Claim 4、\texttt{gemini-3.1-pro} $\to$ Claim 5、\texttt{gemini-3.8-flash} $\to$ Claim 6。large: 同じ順で \texttt{gpt-6-sol-large} $\to$ Claim 7 から \texttt{gemini-3.8-flash-large} $\to$ Claim 12。

% ===== airas: post-experiment sections, filled once runs exist =====

\section{結果}
\label{sec:results}
実験は未実施である。実行後にこの節へ、small の各 run の平均 STE、論文最良行との差、RMS、有効提出件数を \texttt{table1}、large を \texttt{large} として、論文 Table 1 の 6 行と並べた対比表とともに記入する。

\section{考察}
\label{sec:discussion}
実験は未実施である。

\section{結論}
\label{sec:conclusion}
実験は未実施である。

\section*{データの利用可能性}
記録は \airasrecordlink{record.json} にある。実験コード、環境の固定（Dockerfile と \texttt{uv.lock}）、各 run の結果と provenance は同じリポジトリに含まれる。

\bibliographystyle{iclr2024_conference}
\bibliography{references}

\end{document}
